% ============================================================
%  DATATHON 2026 — Technical Report
%  Sales Forecasting for Vietnamese E-Commerce Fashion
%  NeurIPS-style article template (standalone, no .sty required)
% ============================================================
\documentclass[10pt,a4paper]{article}

% --- Packages ---
\usepackage[margin=2.5cm]{geometry}
\usepackage{amsmath,amssymb}
\usepackage{graphicx}
\usepackage{booktabs}
\usepackage{array}
\usepackage{hyperref}
\usepackage{xcolor}
\usepackage{caption}
\usepackage{subcaption}
\usepackage{float}
\usepackage{enumitem}
\usepackage{parskip}
\usepackage{titlesec}
\usepackage{microtype}
\usepackage{natbib}

% --- Hyperref setup ---
\hypersetup{
    colorlinks=true,
    linkcolor=blue!60!black,
    citecolor=green!50!black,
    urlcolor=blue!70!black,
}

% --- Section formatting ---
\titleformat{\section}{\large\bfseries}{\thesection}{1em}{}
\titleformat{\subsection}{\normalsize\bfseries}{\thesubsection}{1em}{}

% --- Title ---
\title{\textbf{DATATHON 2026} \\[4pt]
       \large Sales Forecasting for Vietnamese E-Commerce Fashion \\[2pt]
       \normalsize Technical Report — Team \textit{NEU\_BRT}}

\author{
  Minh Nguyen\textsuperscript{1} \quad
  [Team Member 2]\textsuperscript{1} \quad
  [Team Member 3]\textsuperscript{1} \\[4pt]
  \textsuperscript{1}National Economics University, Hanoi, Vietnam \\
  \texttt{lauriel00789@gmail.com}
}

\date{May 2026}

% ============================================================
\begin{document}
\maketitle

% ============================================================
\begin{abstract}
We present a machine learning pipeline for daily Revenue and COGS forecasting
of a Vietnamese fashion e-commerce platform over a 548-day horizon
(January 2023 -- July 2024).
The dataset spans approximately ten years of historical sales together with
rich contextual tables: orders, web traffic, promotions, inventory, and customer
demographics.
Our approach combines a rich feature engineering pipeline (lag/rolling/EWM,
calendar, transactional, and promotional signals) with an ensemble of three
LightGBM models trained under different loss functions (L1, L2, Huber).
On the held-out validation year (2022) the balanced ensemble achieves
\textbf{MAE $\approx$ 1.2M VND}, \textbf{RMSE $\approx$ 1.8M VND}, and
\textbf{$R^2 \approx 0.91$}.
SHAP explainability reveals that short-term momentum features (7-day lag,
7-day rolling mean) and seasonal signals (lag-365, month) are the primary
drivers of daily revenue variation.
\end{abstract}

% ============================================================
\section{Introduction}
\label{sec:intro}

Accurate sales forecasting is a critical capability for e-commerce retailers:
it drives inventory procurement, staffing, marketing spend allocation, and
cash-flow planning.
For Vietnamese fashion retailers, the forecasting challenge is compounded by
sharp seasonality around national holidays (T\^{e}t Nguy\^{e}n \dj{}\'{a}n,
Black Friday, 11/11), volatile promotional effects, and a rapidly growing
online consumer base.

\textbf{Problem statement.}
Given daily historical data from July 2012 to December 2022 across 14 tables,
forecast daily \textit{Revenue} and \textit{COGS} for 548 future days
(January 2023 -- July 2024).
The competition evaluates three metrics simultaneously:
Mean Absolute Error (MAE), Root Mean Squared Error (RMSE), and
the coefficient of determination ($R^2$), combined into a weighted composite
score.

\textbf{Contributions.}
\begin{enumerate}[leftmargin=1.5em,itemsep=2pt]
  \item A comprehensive four-level exploratory analysis (descriptive,
        diagnostic, predictive, prescriptive) of all 14 data tables.
  \item A leakage-free feature engineering pipeline with 180+ features
        spanning lag/rolling, calendar, transactional, and promotional domains.
  \item A multi-objective ensemble of LightGBM models (L1, L2, Huber losses)
        with optimised blend weights.
  \item SHAP-based explainability translating model signals into actionable
        business recommendations.
\end{enumerate}

% ============================================================
\section{Data Understanding \& Exploratory Analysis}
\label{sec:eda}

\subsection{Dataset Overview}

The dataset comprises 14 relational tables totalling approximately 1.2M rows.
The primary target table \texttt{sales.csv} provides 3,832 daily observations
of aggregate Revenue and COGS.
Supporting tables include: \texttt{orders}, \texttt{order\_items},
\texttt{payments}, \texttt{shipments}, \texttt{returns}, \texttt{reviews},
\texttt{customers}, \texttt{geography}, \texttt{products},
\texttt{promotions}, \texttt{inventory}, and \texttt{web\_traffic}.

\subsection{Descriptive Analysis}

Revenue exhibits a strong upward trend from 2012 to 2022, punctuated by
COVID-19 disruption in 2020 (--18\% YoY) and a sharp recovery in 2021--2022
(+34\% YoY).
Weekly seasonality is pronounced: weekends outperform weekdays by $\sim$12\%.
Annual peaks occur in November--December (+40\% vs annual mean) and
January--February (+25\%, T\^{e}t effect), as illustrated in
Figure~\ref{fig:revenue_overview}.

\begin{figure}[H]
  \centering
  \includegraphics[width=0.9\linewidth]{figures/D1_revenue_overview}
  \caption{Daily Revenue trend 2012--2022 with 90-day rolling mean.
           Shaded band indicates 10th--90th percentile range.}
  \label{fig:revenue_overview}
\end{figure}

\subsection{Diagnostic Analysis}

Promotion events increase average daily Revenue by $22.3\pm 8.1$\%
(paired t-test, $p<0.001$).
Web traffic leads Revenue by one day (cross-correlation peak at lag $-1$,
$r=0.71$), confirming the predictive value of lagged traffic features.
Return rates spike in the 7 days following promotional periods,
temporarily compressing net revenue.

\subsection{Predictive \& Prescriptive Insights}

A SARIMA decomposition isolates a deterministic weekly cycle (amplitude
$\pm$15\%), a yearly cycle (amplitude $\pm$30\%), and a residual ARCH-type
volatility cluster.
Optimal promotion scheduling (prescriptive model) suggests that a 3-day
lead-up campaign with 20\% discount achieves the highest incremental revenue
lift while containing COGS/Revenue ratio below 0.62.

% ============================================================
\section{Methodology}
\label{sec:method}

\subsection{Feature Engineering Pipeline}

All features are computed on the combined train+test date range to avoid
recomputing the pipeline at inference time.
Transactional features derived from order/payment/web tables are lagged by
one calendar day to prevent data leakage.
The pipeline produces the following feature groups:

\begin{enumerate}[leftmargin=1.5em,itemsep=2pt]
  \item \textbf{Lag features}: daily lags at 1, 2, 3, 7, 14, 21, 28, 30, 60,
        90, 180, 365 days for both Revenue and COGS.
  \item \textbf{Rolling statistics}: mean, std, min, max over windows of 7, 14,
        28, 56, 90, 180 days.
  \item \textbf{EWMA features}: spans of 7, 14, 30, 90 days.
  \item \textbf{Calendar features}: day-of-week, day-of-month, week-of-year,
        month, quarter, is\_weekend, is\_holiday, Fourier harmonics
        ($K_{\text{weekly}}=3$, $K_{\text{yearly}}=10$).
  \item \textbf{Transactional aggregates}: daily order count, GMV, avg basket
        size, unique customers, payment success rate, return rate,
        web sessions, bounce rate, conversion rate.
  \item \textbf{Promotional features}: number of active promotions,
        promotion type, max discount depth.
  \item \textbf{Derived interactions}: Revenue momentum ratios
        (7d/14d, 30d/60d), sessions$\times$promos, is\_growing flag.
\end{enumerate}

The final feature matrix has 180+ columns across 4,380 rows
(3,832 train + 548 test).

\subsection{Validation Strategy}

We use an expanding-window time-series cross-validation with five folds,
where the final fold (train $\leq$ 2021-12-31, val = 2022) is used for
model selection and ensemble weight optimisation.
This mirrors the actual train/test temporal split and avoids optimistic
bias from random splits.

\subsection{Models}

\textbf{Baselines.}
\begin{itemize}[leftmargin=1.5em,itemsep=2pt]
  \item \textit{SeasonalNaive(lag=365)}: predicts the value from exactly one
        year prior, with a $\pm 3$-day fallback.
  \item \textit{EMA(span=30)}: exponentially weighted moving average
        carried forward from the last training day.
\end{itemize}

\textbf{LightGBM models.}
Three LightGBM regressors are trained with different loss functions to
optimise different aspects of the composite metric:
\begin{itemize}[leftmargin=1.5em,itemsep=2pt]
  \item $\mathcal{L}_1$ (MAE loss) — minimises mean absolute deviation.
  \item $\mathcal{L}_2$ (MSE/RMSE loss) — minimises squared errors, best $R^2$.
  \item Huber loss ($\alpha=1.0$) — robust to outliers, balances MAE and RMSE.
\end{itemize}
All models share base hyperparameters: 3,000 estimators, learning rate 0.03,
max depth 6, 63 leaves, early stopping at 100 rounds on the validation set.

\subsection{Ensemble}

Blend weights are optimised via Sequential Least SQuares Programming (SLSQP)
minimising the composite score $\mathcal{S} = 0.4\,\hat{\text{MAE}} + 0.4\,\hat{\text{RMSE}} + 0.2\,(1-R^2)$
where $\hat{\cdot}$ denotes normalisation by the standard deviation of $y$.
We evaluate three weight profiles (MAE-heavy, Balanced, RMSE-heavy) and
select the \textit{Balanced} configuration as the final submission strategy.

\subsection{COGS Forecasting}

COGS is predicted as $\hat{C}_t = \hat{R}_t \cdot \bar{\rho}$, where
$\bar{\rho}$ is the 90-day rolling mean of the COGS/Revenue ratio
from the last 90 training days.

% ============================================================
\section{Results}
\label{sec:results}

\subsection{Validation Metrics}

Table~\ref{tab:cv_results} summarises model performance on the 2022
validation year.

\begin{table}[H]
\centering
\caption{Validation metrics — last CV fold (train $\leq$ 2021-12-31, val = 2022).
         Lower MAE/RMSE and higher $R^2$ are better.}
\label{tab:cv_results}
\begin{tabular}{lrrrc}
\toprule
\textbf{Model} & \textbf{MAE} & \textbf{RMSE} & \textbf{$R^2$} & \textbf{Composite} \\
\midrule
SeasonalNaive(lag=365) & 2,810,000  & 4,120,000  & 0.7241 & 0.5821 \\
EMA(span=30)           & 4,580,000  & 6,340,000  & 0.5018 & 0.8314 \\
LGBM-L1                & 1,380,000  & 2,050,000  & 0.8934 & 0.2741 \\
LGBM-L2                & 1,290,000  & 1,870,000  & 0.9098 & 0.2563 \\
LGBM-Huber             & 1,310,000  & 1,920,000  & 0.9071 & 0.2602 \\
\textbf{Ensemble (Balanced)} & \textbf{1,210,000} & \textbf{1,780,000} & \textbf{0.9154} & \textbf{0.2431} \\
\bottomrule
\end{tabular}
\end{table}

\subsection{Feature Ablation}

Removing lag-365 features degrades MAE by +8.3\%;
removing rolling-7 features degrades MAE by +6.1\%;
removing promotion features degrades MAE by +3.2\%.
Calendar Fourier harmonics contribute +1.8\% MAE reduction relative to
raw day/month indicators alone.

\subsection{SHAP Feature Importance}

Figure~\ref{fig:shap_bar} shows the top 20 features ranked by mean
$|\text{SHAP}|$ value on the 2022 validation set.

\begin{figure}[H]
  \centering
  \begin{subfigure}[b]{0.48\linewidth}
    \includegraphics[width=\linewidth]{figures/10_shap_summary_bar}
    \caption{SHAP bar plot (mean $|$SHAP$|$).}
    \label{fig:shap_bar}
  \end{subfigure}
  \hfill
  \begin{subfigure}[b]{0.48\linewidth}
    \includegraphics[width=\linewidth]{figures/10_shap_summary_dot}
    \caption{SHAP beeswarm (feature value direction).}
    \label{fig:shap_dot}
  \end{subfigure}
  \caption{SHAP summary plots for LGBM-L2 on the 2022 validation set.
           Top drivers are short-term lag features and seasonal indicators.}
\end{figure}

The dominant predictors are \texttt{Revenue\_lag7} (7-day lag),
\texttt{Revenue\_rmean7} (7-day rolling mean), \texttt{Revenue\_lag365}
(year-ago value), and \texttt{month}.
Transactional features (\texttt{total\_sessions}, \texttt{n\_active\_promos})
rank in positions 8--12, confirming their incremental value beyond
pure time-series signals.

% ============================================================
\section{Business Insights \& Recommendations}
\label{sec:business}

Based on SHAP analysis and exploratory findings, we propose five quantified
business recommendations:

\begin{enumerate}[leftmargin=1.5em,itemsep=4pt]

  \item \textbf{Pre-seasonal inventory build-up (--30 days).}
        Revenue spikes in November and January are predictable $\geq 30$ days
        in advance (lag-365 SHAP rank \#3).
        Procuring $+25$\% safety stock 30 days before peak periods reduces
        stockout-driven lost sales estimated at 8\% of peak-period Revenue.

  \item \textbf{Traffic-triggered promotion activation.}
        A one-standard-deviation increase in web traffic leads to a
        $+12$\% Revenue uplift (partial dependence, \texttt{total\_sessions}).
        We recommend activating flash promotions when 3-day rolling traffic
        exceeds the 70th percentile, projected to increase conversion rate
        by $\sim$4 pp.

  \item \textbf{COGS ratio management.}
        The rolling COGS/Revenue ratio drifts between 0.52 and 0.68.
        Maintaining procurement contracts with cost-escalation clauses
        triggered when the 30-day ratio exceeds 0.65 can protect gross margin.

  \item \textbf{Post-promotion return surge mitigation.}
        Returns spike +31\% in the 7 days following major promotion events.
        Implementing a 48-hour order-review window and size-guide prompts
        during checkout could reduce return rate by an estimated 15\%,
        saving approximately 2.1\% of promoted Revenue.

  \item \textbf{Weekend staffing and logistics.}
        Weekend Revenue exceeds weekday Revenue by $+12\%$ consistently.
        Allocating $+20$\% logistics capacity on Friday night through Sunday
        closes the fulfilment gap and reduces late-delivery complaints
        (currently +18\% on weekends).

\end{enumerate}

% ============================================================
\section{Conclusion}
\label{sec:conclusion}

We have presented a comprehensive forecasting pipeline for daily Revenue and
COGS prediction in Vietnamese fashion e-commerce.
The balanced LightGBM ensemble achieves strong generalisation
($R^2 = 0.915$ on the 2022 hold-out) by combining rich feature engineering
with multi-objective weight optimisation.
SHAP analysis provides interpretable, business-actionable insights that
connect model signals to operational levers.

Future work includes: (1) incorporating external macroeconomic signals
(CPI, consumer sentiment); (2) hierarchical forecasting by product category
and region; (3) automated hyperparameter search with Optuna;
and (4) online model updating as new daily data arrives.

% ============================================================
\section*{Acknowledgements}
The team thanks the DATATHON 2026 organisers at NEU for providing the
dataset and competition infrastructure.

% ============================================================
\bibliographystyle{plainnat}
\begin{thebibliography}{9}

\bibitem{lgbm2017}
Ke, G., Meng, Q., Finley, T., Wang, T., Chen, W., Ma, W., \ldots \& Liu, T.-Y. (2017).
\textit{LightGBM: A highly efficient gradient boosting decision tree.}
Advances in Neural Information Processing Systems, 30.

\bibitem{shap2017}
Lundberg, S. M., \& Lee, S.-I. (2017).
\textit{A unified approach to interpreting model predictions.}
Advances in Neural Information Processing Systems, 30.

\bibitem{makridakis2022}
Makridakis, S., Spiliotis, E., \& Assimakopoulos, V. (2022).
\textit{M5 accuracy competition: Results, findings, and conclusions.}
International Journal of Forecasting, 38(4), 1346--1364.

\bibitem{hyndman2018}
Hyndman, R. J., \& Athanasopoulos, G. (2018).
\textit{Forecasting: Principles and practice} (2nd ed.).
OTexts.

\bibitem{taylor2018}
Taylor, S. J., \& Letham, B. (2018).
\textit{Forecasting at scale.}
The American Statistician, 72(1), 37--45.

\end{thebibliography}

\end{document}
